\begin{algorithm}[tbp]
\caption{Forward computation and training supervision}
\label{alg:a1}
\small
\textbf{Input:} available sensor observations, model variant, and optional training GT.
\begin{enumerate}[leftmargin=2.4em,label=\arabic*.,itemsep=2pt]
\item Encode X and project BEV features into spatially corresponding patch tokens t.
\item Compute shared c and modality-specific u for every available patch and modality.
\item Predict foreground gate g and modality contributions from tokens and their shared/specific representations.
\item If the variant uses object context, predict context z from its context head.
\item Construct controlled tokens and queries using the main-text forward equations.
\item Apply local cross-modal attention, the enabled context output residual, and PFT.
\item Reassemble the BEV feature and obtain predictions from the detection head.
\item If training:
\item Compute the detector's training losses using GT.
\item If SCL is enabled, add detection losses on the six controlled-token subsets.
\item Construct foreground/context targets from valid GT boxes using Equation~\eqref{eq:a1}.
\item If valid foreground exists and at least two modalities are available:
\item Add foreground representation losses, gate loss, and enabled context loss.
\item Return predictions and the total loss with Table~\ref{tab:a2} coefficients.
\item Return predictions.
\end{enumerate}
The pseudocode specifies data dependencies rather than low-level operator scheduling. Steps 1--7 are shared by training and inference; GT enters supervision from step 9 onward. The target foreground set in step 11 does not replace the predicted gate in step 3. Step 10 follows the SCL implementation in Section~\ref{subsec:a3} and is skipped on V2X. Losses follow main-text Equation~\eqref{eq:7}.
\end{algorithm}
